\documentclass[11pt, a4paper]{article}

% ─── Encodage et langue ───
\usepackage[utf8]{inputenc}
\usepackage[T1]{fontenc}
\usepackage[french]{babel}

% ─── Maths ───
\usepackage{amsmath, amssymb, amsthm}

% ─── Mise en page ───
\usepackage[margin=2.5cm]{geometry}
\usepackage{enumitem}
\usepackage{booktabs}
\usepackage{multirow}
\usepackage{float}
\usepackage{xcolor}
\usepackage{colortbl}
\usepackage{tikz}
\usetikzlibrary{arrows.meta, positioning, fit, backgrounds, calc, decorations.pathreplacing}
\usepackage{pifont}
\usepackage{hyperref}
\hypersetup{colorlinks=true, linkcolor=blue!60!black, urlcolor=blue!70!black}
\usepackage{fancyhdr}

% ─── En-tête / pied de page ───
\pagestyle{fancy}
\fancyhf{}
\fancyhead[L]{\small\textsc{Rapport de synthèse}}
\fancyhead[R]{\small\textsc{Mars 2026}}
\fancyfoot[C]{\thepage}

% ─── Commandes ───
\newcommand{\R}{\mathbb{R}}
\newcommand{\Ltwo}{L^2(\mathcal{T})}
\newcommand{\Dp}{D_p}
\newcommand{\Dw}{D_w}
\newcommand{\DK}{D_K}

% ─── Couleurs ───
\definecolor{func}{HTML}{2166AC}
\definecolor{vect}{HTML}{B2182B}
\definecolor{mix}{HTML}{4DAF4A}
\definecolor{gris}{HTML}{666666}
\definecolor{emphbox}{HTML}{FFF3CD}
\definecolor{emphborder}{HTML}{856404}
\definecolor{resultbox}{HTML}{D4EDDA}
\definecolor{resultborder}{HTML}{155724}

% ─── Boîtes d'emphase ───
\usepackage{framed}

% keybox : fond jaune, cadre ambre
\newenvironment{keybox}[1][]{%
  \def\FrameCommand{\fboxsep=8pt \fcolorbox{emphborder}{emphbox}}%
  \MakeFramed{\advance\hsize-\width \FrameRestore}%
  \noindent\textbf{#1}\par\smallskip\noindent\ignorespaces
}{\endMakeFramed}

% resultbox : fond vert, cadre vert foncé
\newenvironment{resultbox}[1][]{%
  \def\FrameCommand{\fboxsep=8pt \fcolorbox{resultborder}{resultbox}}%
  \MakeFramed{\advance\hsize-\width \FrameRestore}%
  \noindent\textbf{#1}\par\smallskip\noindent\ignorespaces
}{\endMakeFramed}

% ============================================================================
\title{%
  \vspace{-1cm}
  \textbf{Classification non supervisée} \\[0.3em]
  \textbf{de données mixtes fonctionnelles + vectorielles} \\[0.5em]
  \Large Distances hiérarchiques pour données mixtes : \\
  formulation, validation empirique et discussion critique
}
\author{Pierre Chambet}
\date{Mars 2026}

\begin{document}
\maketitle
\thispagestyle{fancy}

\begin{abstract}
\noindent
\textbf{Contexte.} La classification non supervisée de données \textbf{mixtes}
(trajectoires fonctionnelles complétées par des variables vectorielles)
requiert une mesure de dissimilarité adaptée. Ce rapport présente une
construction hiérarchique de distances, puis en évalue les propriétés sur
données réelles et sur un jeu simulé.

\textbf{Résultats principaux.} Sur trois bases étiquetées, la représentation
concaténant scores fonctionnels et variables (stratégie~A) obtient souvent le
meilleur accord avec les classes (ARI), lorsque le nombre de groupes est fixé.
Lorsque les paramètres de distance $(\alpha, \omega)$ sont choisis par
\emph{silhouette}, ils diffèrent systématiquement des couples optimisant l'ARI,
ce qui traduit un décalage entre critère de validation interne et critère
externe. Sous optimisation de l'ARI, le poids de la composante \emph{forme}
($\alpha_{\text{ARI-optimal}} \geq 0.6$ sur les trois jeux) confirme le rôle
informatif des dérivées. Un complément sur \textbf{données simulées}
($p = 20$, répétitions indépendantes) confronte les mêmes familles de méthodes
à un générateur contrôlé ; l'ARI n'y intervient qu'en évaluation, pas en
sélection d'hyperparamètres.

\end{abstract}
\vspace{0.75em}

\tableofcontents
\newpage

\section{Introduction}

\subsection{Données mixtes fonctionnelles et vectorielles}
Dans de nombreux domaines appliqués, un individu statistique est décrit à la
fois par une \textbf{trajectoire} (signal fonctionnel) et par un
\textbf{vecteur de covariables}. Deux formalismes coexistent :
\begin{itemize}
    \item \textbf{Composante fonctionnelle ($Y$)} : processus à temps continu,
    modélisé dans un espace de dimension infinie (ex.\ $L^2$).
    \item \textbf{Composante vectorielle ($X$)} : variables scalaires ou
    vecteur $X \in \mathbb{R}^p$ (covariables cliniques, géographiques, etc.).
\end{itemize}
Analyser séparément la partie fonctionnelle et la partie vectorielle revient
à ignorer d'éventuelles dépendances entre modalités ; l'objectif est d'en
rendre compte \textit{conjointement}.

\subsection{L'approche RS-PCA (référence)}
La littérature classique (Ramsay et Silverman) propose de représenter chaque
courbe par un vecteur de coefficients (bases de B-splines, Fourier, etc.),
puis de concaténer ces coefficients aux variables $X$ et d'appliquer une ACP
euclidienne. La procédure comporte typiquement deux étapes :
\begin{enumerate}
    \item \textbf{Pré-lissage} : projection de la courbe sur une base de
    dimension finie ;
    \item \textbf{Concaténation} : empilement avec $X$ et analyse par ACP
    standard sur le vecteur ainsi formé.
\end{enumerate}

\subsection{Limites pratiques}
Cette voie reste sensible (i) au \textbf{choix de la base} de représentation,
qui conditionne l'ensemble de l'analyse, et (ii) aux données \textbf{sparse
ou irrégulièrement échantillonnées} (SIR) : l'estimation des coefficients peut
devenir instable lorsque l'information le long du temps est incomplète. Un
cadre alternatif consiste à travailler dans un espace hybride sans réduire la
courbe à un simple vecteur de coefficients avant définition de la géométrie.

\section{Cadre théorique : espace hybride $\mathbb{H}$}

\subsection{Objet statistique et espace de Hilbert}
L'Analyse en Composantes Principales hybride fonctionnelle--vectorielle
(HFV-PCA) définit l'individu comme une paire $Z = (Y, X)$ appartenant à un
espace de Hilbert $\mathbb{H}$ combinant la composante fonctionnelle ($L^2$)
et la composante vectorielle ($\mathbb{R}^p$). Cette construction évite de
réduire arbitrairement la courbe avant d'avoir fixé la géométrie de
l'analyse.

\subsection{Produit scalaire hybride}
L'étude de l'inertie (variance) requiert un produit scalaire sur $\mathbb{H}$.
On combine la partie fonctionnelle (intégrale sur le temps) et la partie
vectorielle (produit scalaire euclidien) selon :

\begin{equation}
\langle h_1, h_2 \rangle_{\mathbb{H}} = \underbrace{\langle f_1, f_2 \rangle_{\mathbb{F}}}_{\text{Monde continu (Intégrale)}} + \omega \underbrace{\langle v_1, v_2 \rangle}_{\text{Monde discret (Somme)}}
\end{equation}

\subsection{Pondération $\omega$ entre modalités}
Les échelles de variation du fonctionnel et du vectoriel diffèrent en général ;
sans renormalisation, la modalité de plus forte variance dominerait
numériquement l'autre. Le paramètre $\omega$ calibre le poids relatif
de la partie vectorielle dans le produit scalaire. Une estimation empirique
courante est le ratio des énergies (variations totales) :

\begin{equation}
\hat{\omega} = \frac{\sum_{i=1}^n ||Y_i - \bar{Y}||_{\mathbb{F}}^2}{\sum_{i=1}^n ||X_i - \bar{X}||_{\mathbb{R}^p}^2}
\end{equation}

Ainsi, aucune source n'est éclipsée par la seule unité de mesure.

\subsection*{Notations ($\hat{\cdot}$ et $\tilde{\cdot}$)}
\begin{itemize}
    \item \textbf{Sans accent (ex.\ $\xi_m$)} : quantité théorique (population).
    \item \textbf{Tilde (ex.\ $\tilde{\xi}_m$)} : approximation après troncature
    du nombre de composantes.
    \item \textbf{Chapeau (ex.\ $\hat{\xi}_m$)} : estimateur à partir de
    l'échantillon observé.
\end{itemize}


\section{L'estimation matricielle : du modèle au calcul fini}

Le modèle de l'espace $\mathbb{H}$ est défini en dimension infinie ; en
pratique, on le ramène à un calcul matriciel de dimension finie (réduction
préalable des courbes et des vecteurs, puis covariance sur les résumés). Le
cadre théorique (Jang) justifie cette approximation.

\subsection{Étape 1 : réduction séparée}
\begin{enumerate}
    \item \textbf{Fonctionnel (FPCA)} : extraction des $L$ premiers scores
    $\hat{\eta}_{i}$.
    \item \textbf{Vectoriel (ACP)} : extraction des $J$ premiers scores
    $\hat{\gamma}_{i}$.
\end{enumerate}

\subsection{Étape 2 : vecteur d'état joint}
Les résumés sont empilés en un vecteur de dimension $(L+J)$ :
\begin{equation}
\hat{x}_i = [\underbrace{\hat{\eta}_{i1}, \dots, \hat{\eta}_{iL}}_{\text{Résumés de la courbe}}, \underbrace{\hat{\gamma}_{i1}, \dots, \hat{\gamma}_{iJ}}_{\text{Résumés cliniques}}]^T
\end{equation}

\subsection{Étape 3 : matrice de covariance par blocs}
On estime $\hat{V} = \frac{1}{n-1} \sum_{i=1}^n \hat{x}_i \hat{x}_i^T$, soit la
covariance des coordonnées fonctionnelles et vectorielles réduites. Elle se
partitionne en blocs :
\begin{equation}
V = \begin{bmatrix} V_y & V_{yx} \\ V_{xy} & V_x \end{bmatrix}
\end{equation}

\begin{itemize}
    \item \textbf{Blocs diagonaux ($V_y$, $V_x$)} : variabilité propre au
    fonctionnel et au vectoriel.
    \item \textbf{Hors de la diagonale ($V_{yx}$ et $V_{xy}$)} : covariance
    croisée fonctionnel--vectorielle ; si deux aspects varient ensemble dans
    les données, ces blocs en tiennent compte.
\end{itemize}

\section{Axes principaux et diagonalisation}

Une fois $V$ estimée (y compris les covariances croisées), les directions de
variance maximale sous contrainte $\|e\|=1$ sont données par les
\textbf{vecteurs propres} de $V$ (théorème spectral). Pour une projection sur
$e$, la variance s'écrit $e^T V e$ ; la suite d'axes orthogonaux maximisant
l'inertie expliquée correspond donc à la résolution du problème aux valeurs
propres de $V$.

\subsection{Covariance croisée et axes dominants}
La présence des blocs $V_{yx}$ implique que les directions principales
synthétisent conjointement la variabilité fonctionnelle et vectorielle.

\subsection{Structure d'un vecteur propre}
Chaque vecteur propre $\hat{e}_m$ se décompose selon les blocs fonctionnel et
vectoriel :
\begin{equation}
\hat{e}_m = \begin{bmatrix} \hat{c}_m \\ \hat{d}_m \end{bmatrix}
\end{equation}
\begin{itemize}
    \item \textbf{$\hat{c}_m$} : coefficients sur les fonctions de base de la
    FPCA (contribution fonctionnelle à l'axe).
    \item \textbf{$\hat{d}_m$} : coefficients sur les directions vectorielles
    réduites (contribution des covariables).
\end{itemize}

\section{Reconstruction et scores}

\subsection{Représentation des axes}
Les vecteurs propres dans l'espace réduit se réinterprètent en termes de
fonctions et de variables d'origine.

\subsection{Fonction représentative et profil vectoriel ($\hat{\xi}_m$)}
\begin{enumerate}
    \item \textbf{Composante fonctionnelle ($\hat{\psi}_m(t)$)} : combinaison
    linéaire des fonctions de base FPCA,
    \begin{equation}
    \hat{\psi}_m(t) = \sum_{h=1}^L \hat{c}_{mh} \hat{\phi}_h(t)
    \end{equation}
    \item \textbf{Composante vectorielle ($\hat{\theta}_m$)} : combinaison
    des axes vectoriels réduits.
    \begin{equation}
    \hat{\theta}_m = \sum_{j=1}^J \hat{d}_{mj} \hat{w}_j
    \end{equation}
\end{enumerate}

Ici $\hat{\phi}_h(t)$ et $\hat{w}_j$ désignent respectivement les fonctions de
base issues de la FPCA et les directions vectorielles issues de l'ACP sur
$X$.

\subsection{Scores individuels ($\hat{\rho}_{im}$)}
Le score de l'individu $i$ sur l'axe $m$ est la projection sur $\hat{e}_m$ :
\begin{equation}
\hat{\rho}_{im} = \sum_{h=1}^L \hat{\eta}_{ih} \hat{c}_{mh} + \sum_{j=1}^J \hat{\gamma}_{ij} \hat{d}_{mj}
\end{equation}

% ============================================================================
\section{Construction de distances pour le clustering}
% ============================================================================

Chaque individu $i = 1, \ldots, n$ est décrit par deux objets géométriques
non comparables \emph{a priori} :

\[
  \text{Individu } i : \quad
  \underbrace{\hat{X}_i(t) \in \Ltwo}_{\text{courbe fonctionnelle}}
  \quad + \quad
  \underbrace{Z_i = (z_{i1}, \ldots, z_{ip}) \in \R^p}_{\text{vecteur classique}}.
\]

\begin{figure}[H]
\centering
\begin{tikzpicture}[
    box/.style={draw, rounded corners=6pt, minimum width=4.5cm, minimum height=3cm,
                align=center, font=\small},
  ]
  \node[box, fill=func!8] (ind1) {
    \textbf{Individu $i$} \\[6pt]
    Courbe $X_i(t)$ \\
    {\footnotesize vit dans $L^2$} \\[6pt]
    $Z_i = (48.4, -123.3, 1.6)$ \\
    {\footnotesize vit dans $\R^3$}
  };
  \node[box, fill=vect!8, right=2.5cm of ind1] (ind2) {
    \textbf{Individu $j$} \\[6pt]
    Courbe $X_j(t)$ \\
    {\footnotesize vit dans $L^2$} \\[6pt]
    $Z_j = (43.1, -79.4, 2.8)$ \\
    {\footnotesize vit dans $\R^3$}
  };
  \draw[{Stealth[length=4mm]}-{Stealth[length=4mm]}, thick, dashed, gris]
    (ind1) -- (ind2) node[midway, above, font=\large\bfseries] {$d(i,j) = \;?$};
\end{tikzpicture}
\caption{Le problème central : construire une distance entre objets
qui vivent dans des espaces différents ($L^2$ et $\R^p$).}
\end{figure}

Il n'existe pas de distance euclidienne native entre $L^2$ et $\mathbb{R}^p$.
La suite du rapport définit une \textbf{architecture} de dissimilarités
hiérarchiques :

\begin{enumerate}[nosep]
  \item Mesure la dissimilarité \textbf{entre courbes} (fonctionnel).
  \item Mesure la dissimilarité \textbf{entre vecteurs} (classique).
  \item \textbf{Combine} les deux de manière contrôlée, avec des paramètres
    interprétables.
\end{enumerate}

La solution est construite en \textbf{deux niveaux} hiérarchiques, détaillés
dans les sections suivantes.

% ============================================================================
\newpage
\section{Les briques de base (Niveau 0)}
% ============================================================================

Trois distances élémentaires, chacune capturant un aspect différent des données.

% ────────────────────────────────────────────────────────────────────────
\subsection{$D_0$ : distance $L^2$ entre courbes (\textbf{niveau})}

\[
  D_0(i,j) = \sqrt{\int_{\mathcal{T}} \bigl[\hat{X}_i(t) - \hat{X}_j(t)\bigr]^2 \, dt}
\]

$D_0$ mesure la \textbf{proximité globale des profils} : c'est l'aire entre
les deux courbes (au sens $L^2$). Deux courbes au même niveau mais
légèrement décalées en phase seront proches en $D_0$.

\begin{itemize}[nosep]
  \item[\textcolor{mix}{$\boldsymbol{+}$}] Interprétation directe : ``l'écart vertical moyen''.
  \item[\textcolor{vect}{$\boldsymbol{-}$}] Insensible aux différences de
    \textbf{forme} si les niveaux sont proches.
\end{itemize}

\paragraph{Exemple.}
Vancouver et Montréal ont des formes de courbes similaires (hiver froid, été
chaud) mais Montréal est décalée vers le bas. $D_0$ les sépare bien. En
revanche, Vancouver (courbe plate) et Halifax (forte amplitude) ont des
niveaux proches mais des dynamiques très différentes ; pourtant $D_0$ ne les distingue pas.

% ────────────────────────────────────────────────────────────────────────
\subsection{$D_1$ : distance $L^2$ entre dérivées (\textbf{forme})}

\[
  D_1(i,j) = \sqrt{\int_{\mathcal{T}} \bigl[\hat{X}_i'(t) - \hat{X}_j'(t)\bigr]^2 \, dt}
\]

$D_1$ mesure la proximité des \textbf{dynamiques}. Si $X_i(t)$ est une
température, alors $X_i'(t)$ est la vitesse de changement.

\begin{itemize}[nosep]
  \item[\textcolor{mix}{$\boldsymbol{+}$}] Capture la \textbf{forme},
    indépendamment du niveau vertical.
  \item[\textcolor{vect}{$\boldsymbol{-}$}] Sensible au bruit (la dérivée
    amplifie les fluctuations). Nécessite un bon lissage en amont.
\end{itemize}

\paragraph{Exemple.}
Vancouver (Pacifique) : hivers doux, étés doux $\Rightarrow$ dérivée
\textbf{plate}. Halifax (Atlantique) : hivers froids, étés chauds
$\Rightarrow$ dérivée \textbf{forte} en mars et octobre. $D_0$ les voit
proches (niveaux similaires). $D_1$ les voit \textbf{très différentes}.

% ────────────────────────────────────────────────────────────────────────
\subsection{$D_s$ : distance euclidienne sur $Z$ (\textbf{vectoriel})}

\[
  D_s(i,j) = \left\| Z_i^{\text{std}} - Z_j^{\text{std}} \right\|_2
\]

où $Z^{\text{std}}$ est la standardisation variable par variable
(moyenne~0, variance~1). $D_s$ mesure la proximité dans l'espace des
variables classiques et \textbf{ignore totalement les courbes}. Cette
standardisation met toutes les composantes de $Z$ sur une échelle
comparable avant qu'elles soient combinées au fonctionnel dans les
stratégies mixtes.

% ────────────────────────────────────────────────────────────────────────
\subsection{Bilan : aucune ne suffit seule}

\begin{table}[H]
\centering
\begin{tabular}{lll}
  \toprule
  \textbf{Distance} & \textbf{Ce qu'elle voit} & \textbf{Ce qu'elle rate} \\
  \midrule
  $D_0$ & Niveau vertical des courbes & Forme / dynamique \\
  $D_1$ & Forme / vitesse de changement & Niveau vertical \\
  $D_s$ & Position dans $\R^p$ & Les courbes entièrement \\
  \bottomrule
\end{tabular}
\caption{Complémentarité des trois distances de base.}
\end{table}

Aucune des trois ne capture toute l'information $\Rightarrow$ il faut les
\textbf{combiner}.

% ============================================================================
\newpage
\section{$\Dp(\alpha)$ : combiner niveau et forme (Niveau 1)}
% ============================================================================

\subsection{Le problème d'échelle}

$D_0$ et $D_1$ sont sur des échelles très différentes. Sur Canadian Weather :
\[
  D_0 \in [3.9,\; 532] \qquad D_1 \in [0.51,\; 6.6] \qquad \text{(facteur $\times 80$)},
\]
où le minimum de $D_0$ exclut la diagonale (distance à soi-même nulle).

Si on les additionne brut, $D_1$ ne pèse rien. \textbf{Solution :}
normaliser par le maximum :
\[
  \tilde{D}_0 = \frac{D_0}{\max(D_0)} \in [0,1], \qquad
  \tilde{D}_1 = \frac{D_1}{\max(D_1)} \in [0,1].
\]

\subsection{La formule}

\begin{keybox}[Définition de $\Dp(\alpha)$]
\[
  \Dp(\alpha) = \sqrt{(1 - \alpha)\,\tilde{D}_0^{\,2}
  + \alpha\,\tilde{D}_1^{\,2}},
  \qquad \alpha \in [0, 1]
\]
\end{keybox}

La normalisation de $D_0$ et $D_1$ dans $[0,1]$ avant combinaison
garantit qu'aucune des deux composantes (niveau ou forme) ne domine
mécaniquement la distance. C'est une condition importante pour que les
stratégies mixtes du niveau~2 ($\Dw$ et $\DK$) puissent pondérer
proprement le fonctionnel par rapport au vectoriel.

\subsection{Le curseur $\alpha$}

Le paramètre $\alpha$ agit comme un \textbf{curseur} entre le niveau
et la forme :

\begin{figure}[H]
\centering
\begin{tikzpicture}
  \draw[very thick, func!70] (0,0) -- (10,0);
  \foreach \x/\lab in {0/$\alpha{=}0$, 5/$\alpha{=}0.5$, 10/$\alpha{=}1$} {
    \fill[func] (\x, 0) circle (3pt);
    \node[below=6pt, font=\small] at (\x, 0) {\lab};
  }
  \node[above=8pt, font=\small\bfseries, func] at (0, 0) {Niveau pur};
  \node[above=8pt, font=\small] at (5, 0) {Compromis};
  \node[above=8pt, font=\small\bfseries, func] at (10, 0) {Forme pure};

  \node[below=22pt, font=\footnotesize, gris] at (0, 0) {$\Dp = \tilde{D}_0$};
  \node[below=22pt, font=\footnotesize, gris] at (10, 0) {$\Dp = \tilde{D}_1$};
\end{tikzpicture}
\end{figure}

\subsection{Propriétés}

$\Dp(\alpha)$ est une \textbf{vraie distance} (positivité, symétrie,
inégalité triangulaire par Minkowski sur $\ell^2$ pondéré). C'est important
car PAM nécessite une matrice de distances au sens métrique.

\begin{keybox}[Lien avec la littérature (Nguyen, 2023)]
Une construction analogue de $\Dp$ n'associait que des distances
\textbf{purement fonctionnelles} ($D_0$ et $D_1$), sans partie vectorielle
explicite.
\end{keybox}

% ============================================================================
\newpage
\section{Le mix total : $\Dw$ et $\DK$ (Niveau 2)}
% ============================================================================

On a maintenant $\Dp(\alpha)$ (fonctionnel combiné) et $D_s$ (vectoriel).
Deux stratégies pour les combiner.

\subsection{Architecture complète}

\begin{figure}[H]
\centering
\begin{tikzpicture}[
    base/.style={draw, rounded corners=4pt, minimum width=2.5cm, minimum height=0.9cm,
                 align=center, font=\small\bfseries},
    fleche/.style={-{Stealth[length=3mm]}, thick},
    etiq/.style={font=\footnotesize\itshape, text=gris},
  ]

  \node[base, fill=func!15, text=func] (D0) at (0,0) {$D_0$ \footnotesize niveau};
  \node[base, fill=func!15, text=func] (D1) at (3.5,0) {$D_1$ \footnotesize forme};
  \node[base, fill=vect!15, text=vect] (Ds) at (8.5,0) {$D_s$ \footnotesize vectoriel};

  \node[base, fill=func!25, text=func] (Dp) at (1.75,-1.8) {$\Dp(\alpha)$};
  \node[base, fill=mix!20, text=mix!70!black] (Dw) at (4.0,-3.6) {$\Dw(\alpha, \omega)$};
  \node[base, fill=mix!20, text=mix!70!black] (DK) at (7.5,-3.6) {$\DK(\alpha)$};

  \draw[fleche, func!70] (D0) -- (Dp) node[etiq, midway, left=2pt] {$1{-}\alpha$};
  \draw[fleche, func!70] (D1) -- (Dp) node[etiq, midway, right=2pt] {$\alpha$};
  \draw[fleche, func!70] (Dp) -- (Dw) node[etiq, midway, left=2pt] {$\omega$};
  \draw[fleche, vect!70] (Ds) -- (Dw) node[etiq, midway, right=2pt] {$1{-}\omega$};
  \draw[fleche, func!70] (Dp) -- (DK) node[etiq, midway, left=2pt] {$K_f$};
  \draw[fleche, vect!70] (Ds) -- (DK) node[etiq, midway, right=2pt] {$K_s$};

  \node[font=\footnotesize\scshape, text=gris] at (-1.5,0) {Niv.\ 0};
  \node[font=\footnotesize\scshape, text=gris] at (-1.5,-1.8) {Niv.\ 1};
  \node[font=\footnotesize\scshape, text=gris] at (-1.5,-3.6) {Niv.\ 2};
\end{tikzpicture}
\caption{Architecture à deux niveaux.
\textcolor{func}{Bleu} = fonctionnel,
\textcolor{vect}{rouge} = vectoriel,
\textcolor{mix}{vert} = mixte.}
\end{figure}

On retrouve l'esprit d'une \textbf{RS-PCA} : les courbes sont d'abord résumées
(FPCA), puis confrontées au vecteur $Z$. La différence entre les stratégies
A, B et C est la manière d'\textbf{exploiter} ce résumé : coordonnées
euclidiennes + k-means (A), ou matrices de distances / noyaux construits sur
les mêmes ingrédients (B et C).

% ────────────────────────────────────────────────────────────────────────
\subsection{Résumé des trois stratégies de clustering}

\begin{table}[H]
\centering
\renewcommand{\arraystretch}{1.25}
\begin{tabular}{l p{3.2cm} p{3.4cm} p{3.4cm}}
  \toprule
  \textbf{Stratégie} & \textbf{Représentation} &
  \textbf{Méthode de clustering} & \textbf{Rôle de la FPCA / distances} \\
  \midrule
  A (FPCA+Z) &
  Vecteurs $W = [\xi\_{\text{FPCA}} \mid Z]$ standardisés &
  k-means (distance euclidienne) &
  FPCA sur les courbes, $Z$ standardisé, pas d'ACP sur $Z$ \\
  B ($\Dw$) &
  Matrice de distances $\Dw(\alpha,\omega)$ &
  PAM (k-médoïdes) sur distances &
  Combinaison pondérée de $\Dp(\alpha)$ (fonctionnel) et $D_s$ (vectoriel) \\
  C ($\DK$) &
  Matrice de distances $\DK(\alpha)$ &
  PAM (k-médoïdes) sur distances &
  Produit de noyaux issus de $\Dp(\alpha)$ et $D_s$ puis retour à une distance \\
  \bottomrule
\end{tabular}
\caption{Résumé des trois stratégies : type de représentation, méthode de
clustering et usage de la FPCA ou des distances.}
\end{table}

% ────────────────────────────────────────────────────────────────────────
\subsection{Stratégie A : FPCA + $Z$, concaténation dans $\R^{K+p}$}

La première stratégie mixte correspond à une \textbf{RS-PCA (Reduced-Space
PCA) au sens strict}. La chaîne est la suivante :
\begin{enumerate}[nosep]
  \item lissage des courbes $X_i(t)$ puis FPCA $\Rightarrow$ scores
        fonctionnels $(\xi_{i1}, \ldots, \xi_{iK})$ ;
  \item construction de $Z$ (géographie, spectres agrégés, etc.) et
        standardisation variable par variable ;
  \item concaténation des deux blocs dans un même vecteur ;
  \item standardisation colonne par colonne, puis k-means avec distance
        euclidienne.
\end{enumerate}

Concrètement, avant standardisation finale, on forme pour chaque
individu :
\[
  U_i = (\xi_{i1}, \ldots, \xi_{iK}, Z_i) \in \R^{K+p},
\]
puis on applique k-means sur la matrice $W$ obtenue en concaténant les
scores FPCA et $Z$, tous deux centrés-réduits. Autrement dit, la
distance ``vue'' par la stratégie A est la distance euclidienne dans un
espace réduit où chaque coordonnée (score fonctionnel ou variable de
$Z$) contribue de manière comparable.

Il n'y a \textbf{pas d'ACP supplémentaire sur $Z$} : les variables
vectorielles sont uniquement standardisées, ce qui garde des coordonnées
interprétables lorsque $p$ est modeste. La stratégie A sert de référence
\textbf{mixte} à comparer à B et C, qui combinent les sources au niveau des
distances ou des noyaux.

% ────────────────────────────────────────────────────────────────────────
\subsection{Stratégie B : $\Dw(\alpha, \omega)$, pondération directe}

\begin{keybox}[Définition de $\Dw(\alpha, \omega)$]
\[
  \Dw(\alpha, \omega) = \sqrt{\omega \cdot \Dp(\alpha)^2
  + (1 - \omega) \cdot \tilde{D}_s^{\,2}},
  \qquad \alpha, \omega \in [0, 1]
\]
\end{keybox}

Les deux entrées de cette formule sont donc des \textbf{distances déjà
mise à l'échelle} : $\Dp(\alpha)$, construite à partir de $\tilde{D}_0$
et $\tilde{D}_1$ dans $[0,1]$, et $\tilde{D}_s$, issue de $Z$
standardisé. Le paramètre $\omega$ joue alors le rôle de curseur
explicite entre la contribution fonctionnelle et la contribution
vectorielle, sur une base devenue comparable.

Les paramètres $\alpha$ et $\omega$ jouent des rôles distincts :

\begin{figure}[H]
\centering
\begin{tikzpicture}
  % Bouton alpha
  \draw[very thick, func!70] (0, 0.8) -- (5, 0.8);
  \fill[func] (0, 0.8) circle (2.5pt); \fill[func] (5, 0.8) circle (2.5pt);
  \node[left, font=\small\bfseries] at (-0.3, 0.8) {$\alpha$ :};
  \node[above=3pt, font=\footnotesize] at (0, 0.8) {Niveau};
  \node[above=3pt, font=\footnotesize] at (5, 0.8) {Forme};
  \node[right=3pt, font=\footnotesize, gris] at (5.3, 0.8) {(dans le fonctionnel)};

  % Bouton omega
  \draw[very thick, mix!70] (0, 0) -- (5, 0);
  \fill[mix!70!black] (0, 0) circle (2.5pt); \fill[mix!70!black] (5, 0) circle (2.5pt);
  \node[left, font=\small\bfseries] at (-0.3, 0) {$\omega$ :};
  \node[above=3pt, font=\footnotesize] at (0, 0) {Vecteurs $Z$};
  \node[above=3pt, font=\footnotesize] at (5, 0) {Courbes};
  \node[right=3pt, font=\footnotesize, gris] at (5.3, 0) {(entre les 2 sources)};
\end{tikzpicture}
\end{figure}

Cas limites :

\begin{center}
\begin{tabular}{ccl}
  \toprule
  $\boldsymbol{\alpha}$ & $\boldsymbol{\omega}$ & \textbf{Distance effective} \\
  \midrule
  $*$ & $0$ & $\tilde{D}_s$ seul (courbes ignorées) \\
  $0$ & $1$ & $\tilde{D}_0$ seul (niveau des courbes) \\
  $1$ & $1$ & $\tilde{D}_1$ seul (forme des courbes) \\
  $0.5$ & $0.5$ & Compromis global \\
  \bottomrule
\end{tabular}
\end{center}

\textbf{Sélection des paramètres :} balayage systématique sur une grille
régulière de $(\alpha, \omega)$, en maximisant la \textbf{silhouette moyenne}
(critère non supervisé).

% ────────────────────────────────────────────────────────────────────────
\subsection{Stratégie C : $\DK(\alpha)$, noyaux gaussiens}

Au lieu de pondérer des distances, on passe par des \textbf{similarités}
(noyaux), qu'on combine par produit :

\begin{align}
  K_f(i,j) &= \exp\!\left( -\frac{\Dp(\alpha)_{ij}^2}{2\sigma_f^2} \right),
  & \sigma_f &= \text{médiane}\{\Dp(\alpha)_{ij}\},  \\[4pt]
  K_s(i,j) &= \exp\!\left( -\frac{D_s(i,j)^2}{2\sigma_s^2} \right),
  & \sigma_s &= \text{médiane}\{D_s(i,j)\},  \\[4pt]
  K &= K_f \cdot K_s, & & \text{(produit = ``ET logique flou'')} \\[4pt]
  \DK(i,j) &= \sqrt{K(i,i) + K(j,j) - 2K(i,j)}.
\end{align}

Le choix $\sigma_f = \text{médiane}\{\Dp(\alpha)_{ij}\}$ et
$\sigma_s = \text{médiane}\{D_s(i,j)\}$ joue le rôle de
\textbf{mise à l'échelle automatique} dans chaque noyau : pour chaque
source, des distances typiques correspondent à des similarités de taille
intermédiaire. Le produit $K = K_f \cdot K_s$ signifie ensuite que deux
individus ne sont similaires \textbf{que si les deux sources le
confirment}. Si $K_f \approx 0$ (courbes très différentes), alors
$K \approx 0$ même si $K_s = 1$.

\textbf{Avantage :} pas de paramètre $\omega$ ; balayage sur $\alpha$ seul
(une grille plus courte qu'en B).

\begin{table}[H]
\centering
\begin{tabular}{lcccc}
  \toprule
  & \textbf{Paramètres} & \textbf{Combinaison} & \textbf{Recherche sur grille} \\
  \midrule
  $\Dw$ & $\alpha, \omega$ & Somme pondérée & Grille 2D \\
  $\DK$ & $\alpha$ + $\sigma$ auto & Produit de noyaux & Grille 1D \\
  \bottomrule
\end{tabular}
\caption{Comparaison des deux stratégies de combinaison.}
\end{table}

% ────────────────────────────────────────────────────────────────────────
\subsection{Arbitrages utiles à retenir}

\begin{itemize}[nosep]
  \item Les variables $Z$ restent en coordonnées interprétables (pas d'ACP
    sur $Z$ lorsque $p$ est faible).
  \item La stratégie A utilise des coordonnées euclidiennes dans un espace
    réduit (k-means). Les stratégies B et C s'appuient sur des matrices de
    distances (PAM / k-médoïdes), adaptées lorsque la distance est construite
    directement.
  \item Côté fonctionnel, $D_0$ et $D_1$ sont normalisées avant combinaison ;
    côté vectoriel, $Z$ est standardisé. Les paramètres $\omega$ (B) et les
    échelles des noyaux (C) règlent le poids fonctionnel vs vectoriel.
\end{itemize}

% ============================================================================
\newpage
\section{Résultats sur trois jeux de données}
% ============================================================================

\subsection{Jeux de données et protocole}

Trois jeux publics, disposant d'étiquettes utilisées uniquement pour
l'évaluation (ARI), couvrent des contextes applicatifs variés :

\begin{table}[H]
\centering
\renewcommand{\arraystretch}{1.2}
\begin{tabular}{llcccc}
  \toprule
  \textbf{Dataset} & \textbf{Courbes $X(t)$} & $n$ & \textbf{$Z$} & $k$ & \textbf{Classes} \\
  \midrule
  Canadian Weather & Température (365j) & 35 & lat, lon, précip & 4 & Régions clim. \\
  Growth & Taille enfant (31 âges) & 93 & taille fin., croiss. & 2 & Garçon / Fille \\
  Tecator & Spectres NIR (100$\lambda$) & 215 & water, protein & 3 & Niv.\ de gras \\
  \bottomrule
\end{tabular}
\caption{Les trois jeux de données utilisés pour la validation.}
\end{table}

Pour chaque jeu : lissage des courbes, FPCA, calcul des distances, recherche
des paramètres $(\alpha, \omega)$ pour la stratégie concernée, clustering,
puis évaluation par silhouette (critère non supervisé) et par ARI (accord avec
les vraies classes, à des fins de comparaison seulement).

\subsection{Tableau récapitulatif}

\begin{table}[H]
\centering
\resizebox{\textwidth}{!}{%
\renewcommand{\arraystretch}{1.3}
\begin{tabular}{l l cc cc cc cc cc}
  \toprule
  & & \multicolumn{2}{c}{\textbf{Baseline $D_0$}} & \multicolumn{2}{c}{\textbf{Baseline $D_s$}}
    & \multicolumn{2}{c}{\textbf{A (FPCA+Z)}} & \multicolumn{2}{c}{\textbf{B ($\Dw$)}} & \multicolumn{2}{c}{\textbf{C ($\DK$)}} \\
  \cmidrule(lr){3-4}\cmidrule(lr){5-6}\cmidrule(lr){7-8}\cmidrule(lr){9-10}\cmidrule(lr){11-12}
  \textbf{Dataset} & $n$ / $k$ & Sil. & ARI & Sil. & ARI & Sil. & ARI & Sil. & ARI & Sil. & ARI \\
  \midrule
  Canadian Weather & 35 / 4 & 0.285 & 0.229 & 0.448 & 0.616 & 0.395 & \textbf{0.748} & 0.448 & 0.616 & 0.329 & 0.686 \\
  Growth           & 93 / 2 & 0.402 & 0.115 & 0.554 & 0.424 & 0.360 & \textbf{0.682} & 0.554 & 0.424 & 0.261 & 0.368 \\
  Tecator          & 215 / 3 & 0.509 & 0.151 & 0.476 & 0.546 & 0.403 & \textbf{0.439} & 0.509 & 0.151 & 0.266 & 0.371 \\
  \bottomrule
\end{tabular}%
}

\caption{Résultats sur les trois jeux (tirage aléatoire fixé pour la
reproductibilité). Pour B, $(\alpha, \omega)$ est choisi par silhouette.
Les valeurs \textbf{en gras} : meilleur ARI par jeu. La stratégie A est la
meilleure sur les trois jeux ; pour B, la silhouette ramène souvent
vers les baselines ($\omega \in \{0, 1\}$).}
\label{tab:resultats}
\end{table}

\subsection{Synthèse des résultats}

\paragraph{(i) Distance $D_0$ seule.}
L'ARI associé à $D_0$ reste inférieur à celui d'au moins une autre méthode sur
chaque jeu : la seule information de \emph{niveau} est insuffisante.

\paragraph{(ii) Stratégie A.}
La concaténation FPCA + $Z$ avec k-means obtient le meilleur ARI sur les trois
jeux : Canadian ($0.748$), Growth ($0.682$), Tecator ($0.439$).

\paragraph{(iii) Poids de la forme sous ARI optimal.}
Les $\alpha$ maximisant l'ARI satisfont $\alpha \geq 0.6$ sur les trois jeux.

\begin{figure}[H]
\centering
\begin{tikzpicture}
  \draw[thick, ->] (0,0) -- (11, 0) node[right] {$\alpha$};
  \foreach \x in {0,1,...,10} {
    \draw (\x, -0.1) -- (\x, 0.1);
    \node[below=4pt, font=\tiny] at (\x, -0.1) {\pgfmathparse{\x/10}\pgfmathprintnumber[fixed, precision=1]{\pgfmathresult}};
  }
  % Marques
  \fill[func] (6, 0.4) circle (4pt) node[above=4pt, font=\footnotesize] {CAN};
  \fill[func] (9, -0.5) circle (4pt) node[below=4pt, font=\footnotesize] {GRO};
  \fill[func] (10, 0.4) circle (4pt) node[above=4pt, font=\footnotesize] {TEC};

  % Zone
  \draw[very thick, mix!60, decorate, decoration={brace, amplitude=8pt, mirror}]
    (5.5, -1) -- (10.5, -1) node[midway, below=10pt, font=\small\bfseries, mix!70!black]
    {Zone optimale};
\end{tikzpicture}
\caption{Les $\alpha$ ARI-optimaux tombent tous dans $[0.6, 1.0]$.
Les dérivées ($D_1$) sont \textbf{systématiquement informatives}.}
\end{figure}

\paragraph{(iv) Sélection par silhouette (stratégie B).}
La silhouette conduit à $\omega \in \{0, 1\}$ sur les trois jeux (retour aux
baselines $D_s$ ou $D_0$). La famille $\Dw$ n'est pas en cause en soi : la
difficulté réside dans le \textbf{critère de choix} de $(\alpha, \omega)$.

% ============================================================================
\newpage
\section{Données simulées}
\label{sec:exp3}
% ============================================================================

Les trois jeux réels ne permettent pas de régler finement le niveau de bruit
ni le couplage entre signal fonctionnel et partie vectorielle. Un
\textbf{complément sur données simulées} répond à ce besoin : générateur à
paramètres connus, ici avec $p = 20$ variables, $n = 300$ individus et $K = 3$
classes. Quatre scénarios S1--S4 modulent la force du signal (versions plus
ou moins favorables au fonctionnel, au vectoriel, ou aux deux). Pour chaque
scénario, $50$ répétitions indépendantes sont effectuées ; on applique la
même logique que sur données réelles (lissage, réduction fonctionnelle,
distances, clustering) et l'on compare baselines ($D_0$, $D_1$, $D_s$,
$\mathrm{Df}_{\mathrm{sil}}$), stratégies A, B, C, ainsi qu'une variante
exploratoire combinant réduction sur $Z$ et noyaux (\emph{DK reconstruit} dans
les tableaux).

\begin{keybox}[Protocole d'évaluation]
Sur ce benchmark simulé, les hyperparamètres des grilles ($\alpha$,
$(\alpha,\omega)$ pour B, etc.) sont choisis par \textbf{silhouette},
comme pour les résultats du tableau~\ref{tab:resultats}. L'\textbf{ARI}
sert uniquement à mesurer l'accord avec les classes simulées ; il n'est
\textbf{pas} utilisé pour sélectionner les paramètres.
\end{keybox}

\begin{table}[H]
\centering
\renewcommand{\arraystretch}{1.25}
\begin{tabular}{lcc}
  \toprule
  \textbf{Méthode} & \textbf{ARI moyen (global)} & \textbf{Sil.\ moyen} \\
  \midrule
  C ($\DK$, référence) & 0.651 & 0.126 \\
  $D_1$ (forme seule)               & 0.626 & 0.257 \\
  DK reconstruit (HFV + noyaux)    & 0.608 & 0.118 \\
  A (FPCA + $Z$)                   & 0.561 & 0.134 \\
  B ($\Dw$, sil.)                  & 0.555 & 0.287 \\
  $\mathrm{Df}_{\mathrm{sil}}$     & 0.555 & 0.287 \\
  $D_0$                            & 0.528 & 0.256 \\
  $D_s$                            & 0.414 & 0.116 \\
  \bottomrule
\end{tabular}
\caption{Données simulées : moyennes sur l'ensemble des scénarios et des
répétitions ($4 \times 50$ jeux).}
\label{tab:sim_global}
\end{table}

\begin{table}[H]
\centering
\renewcommand{\arraystretch}{1.2}
\small
\begin{tabular}{lp{7.2cm}}
  \toprule
  \textbf{Scénario} & \textbf{Trois meilleures méthodes (ARI moyen)} \\
  \midrule
  S1 & C ($0.94$) $>$ DK reconstruit ($0.90$) $>$ B / $D_1$ / $\mathrm{Df}_{\mathrm{sil}}$ ($0.87$) \\
  S2 & C ($0.89$) $>$ B / $D_1$ / $\mathrm{Df}_{\mathrm{sil}}$ ($0.87$) $>$ $D_0$ ($0.83$) \\
  S3 & A ($0.62$) $>$ $D_s$ ($0.51$) $>$ C ($0.46$) \\
  S4 & $D_1$ ($0.38$) $>$ C ($0.32$) $>$ A ($0.32$) \\
  \bottomrule
\end{tabular}
\caption{Classement par scénario (ARI moyen sur $50$ répétitions ; arrondis).}
\label{tab:sim_ranking}
\end{table}

Sur S1--S2, la stratégie C et la variante DK reconstruit tirent surtout parti
du signal fonctionnel ; lorsqu'il s'affaiblit (S3--S4), toutes les méthodes
se dégradent et les écarts se resserrent, ce qui est cohérent avec l'importance des
dérivées observée sur données réelles. Les tableaux détaillés et le protocole
complet de cette étude simulée sont documentés à part (hors de ce rapport).

% ============================================================================
\newpage
\section{Discussion : silhouette et calibrage des distances}
% ============================================================================

\subsection{Définition}

Pour chaque point, la silhouette compare la distance intra-cluster à la
distance au cluster le plus proche ; elle mesure une \textbf{compacité et une
séparation géométriques} du partitionnement obtenu.

\subsection{Illustration : Canadian Weather}

\begin{table}[H]
\centering
\renewcommand{\arraystretch}{1.3}
\begin{tabular}{p{5.5cm} p{5.5cm}}
  \toprule
  \textbf{Sil-optimal : $(\alpha{=}0, \omega{=}0)$} &
  \textbf{ARI-optimal : $(\alpha{=}0.6, \omega{=}0.8)$} \\
  \midrule
  Distance $= \tilde{D}_s$ (que du vectoriel) &
  Distance $= 80\%$ fonctionnel + $20\%$ vectoriel \\[4pt]
  Clusters par \textbf{géographie} &
  Clusters par \textbf{dynamique climatique} \\[4pt]
  Géométriquement \textbf{compacts} &
  Géométriquement \textbf{allongés} \\[4pt]
  Silhouette $= 0.448$ \textcolor{mix}{$\boldsymbol{\checkmark}$} &
  Silhouette $= 0.312$ \\[4pt]
  ARI $= 0.616$ &
  ARI $= 0.898$ \textcolor{mix}{$\boldsymbol{\checkmark}$} \\[2pt]
  \small 9 stations mal classées &
  \small \textbf{1 seule} station mal classée \\
  \bottomrule
\end{tabular}
\caption{Canadian Weather : la silhouette préfère des clusters \emph{ronds}
à des clusters \emph{justes}. L'ARI $= 0.898$ est obtenu en maximisant l'ARI
sur la grille $(\alpha, \omega)$, et non par sélection silhouette. Le couple
ARI-optimal classe correctement 34 stations sur 35.}
\end{table}

\begin{keybox}[Synthèse]
La silhouette favorise la \textbf{séparation géométrique} des groupes, pas
nécessairement leur adéquation aux classes de référence. Elle reste pertinente
pour le choix de $k$, mais se révèle peu fiable pour calibrer
$(\alpha, \omega)$.
\end{keybox}

\subsection{Conséquence}

\begin{resultbox}[Conclusion]
L'architecture de distances considérée peut atteindre un ARI élevé
(Canadian Weather : jusqu'à $\approx 0.9$ avec des paramètres adaptés). La
principale difficulté réside dans la \textbf{sélection non supervisée} de
$(\alpha, \omega)$, distincte de celle suggérée par la silhouette.
\end{resultbox}

\end{document}
